\section{Dataset Schema}
\label{app:dataset-schema}

We release the audited benchmark as a Hugging Face dataset. The repository URL
is withheld for anonymous review and provided in the camera-ready version. The
dataset publishes two named splits. The \texttt{test} split holds the
100-paper frozen benchmark and carries \texttt{split="eval"} in each row. The
\texttt{validation} split holds the disjoint 14-paper development set and
carries \texttt{split="dev"}. There is no \texttt{train} split. Each row is
one paper, keyed by its arXiv identifier in \texttt{custom\_id}. The row
records the paper's central claim, the audited artifact-availability signals,
the reproduction task handed to the agent, the pinned match target, and the
compute estimate that placed the paper in its band.

Table~\ref{tab:schema-toplevel} lists the top-level fields. Four of them are
nested objects, expanded in
Tables~\ref{tab:schema-links}--\ref{tab:schema-h100}.

\begin{table}[t]
\caption{Top-level fields of each dataset row.}
\label{tab:schema-toplevel}
\begin{center}
\small
\begin{tabular}{@{}lll@{}}
\toprule
Field & Type & Description \\
\midrule
\texttt{custom\_id} & string & arXiv id, the primary key. \\
\texttt{central\_claim} & string & The paper's central empirical claim in one sentence. \\
\texttt{claim\_evidence} & string & Quoted paper text grounding the claim. \\
\texttt{paper\_kind} & string & Paper type. All rows are \texttt{empirical}. \\
\texttt{mre\_config} & string & Prose recipe for the cheapest configuration that tests the claim. \\
\texttt{agent\_task} & string & The reproduction instruction handed to the agent. \\
\texttt{verified\_links} & object & Tool-confirmed artifact URLs (Table~\ref{tab:schema-links}). \\
\texttt{signals} & object & Per-artifact availability audit (Table~\ref{tab:schema-signals}). \\
\texttt{match\_target} & object & The quantity a run must match (Table~\ref{tab:schema-matchtarget}). \\
\texttt{h100\_estimate} & object & Compute estimate and its basis (Table~\ref{tab:schema-h100}). \\
\texttt{score} & int & Artifact-availability score, 0 (all present) to 3 (most missing). \\
\texttt{tier} & string & \texttt{Easy}, \texttt{Medium}, or \texttt{Hard}, derived from \texttt{score}. \\
\texttt{split} & string & \texttt{eval} (test split) or \texttt{dev} (validation split). \\
\texttt{selection\_band} & string & Compute band used for stratified selection. \\
\texttt{h100\_band} & string & Compute band, one of \texttt{0-8}, \texttt{8-32}, \texttt{32-96}, \texttt{96-192}. \\
\texttt{h100\_hours\_estimate} & float & Estimated H100-hours to run the match target. \\
\texttt{h100\_estimate\_basis} & string & Free-text justification for the estimate. \\
\texttt{audited\_h100\_hours} & float & Human-adjudicated H100-hours. \\
\texttt{h100\_recomputed\_hours} & float or null & Recomputed hours when the estimate was revised. \\
\texttt{h100\_arithmetic\_mismatch} & bool or null & Flag set when estimate arithmetic disagreed. \\
\texttt{h100\_hours\_adjudicated} & bool & Whether a human adjudicated the compute figure. \\
\texttt{h100\_needs\_human\_review} & bool & Whether the estimate was flagged for review. \\
\texttt{verification\_status} & string & Classifier verification state. All rows are \texttt{verified}. \\
\texttt{web\_verification} & string & Web-retrieval state. All rows are \texttt{available}. \\
\texttt{exit\_reason} & string & Classifier exit reason. All rows are \texttt{natural}. \\
\bottomrule
\end{tabular}
\end{center}
\end{table}

\begin{table}[t]
\caption{\texttt{verified\_links}: artifact URLs confirmed by a tool call. Each
field is a list of strings and is empty when no artifact of that kind was
located.}
\label{tab:schema-links}
\begin{center}
\small
\begin{tabular}{@{}ll@{}}
\toprule
Field & Description \\
\midrule
\texttt{paper\_or\_project} & Paper or project-page URLs. \\
\texttt{code} & Code repository URLs. \\
\texttt{dataset} & Dataset URLs. \\
\texttt{weights} & Model-weight URLs. \\
\bottomrule
\end{tabular}
\end{center}
\end{table}

\begin{table}[t]
\caption{\texttt{signals}: the per-artifact availability audit. The four
audited signals are \texttt{code\_available}, \texttt{dataset\_available},
\texttt{weights\_available}, and \texttt{dataset\_is\_standard}. Each signal
is an object with the three fields below. The audit trusts a tool call, not
the paper's own statement, so a signal is true only when its
\texttt{verification} field is \texttt{tool\_verified}.}
\label{tab:schema-signals}
\begin{center}
\small
\begin{tabular}{@{}lll@{}}
\toprule
Field & Type & Description \\
\midrule
\texttt{value} & bool & Whether the artifact is present. \\
\texttt{verification} & string & \texttt{tool\_verified}, \texttt{tool\_searched\_not\_found}, or \texttt{not\_applicable}. \\
\texttt{evidence} & string & The tool observation supporting the value. \\
\bottomrule
\end{tabular}
\end{center}
\end{table}

\begin{table}[t]
\caption{\texttt{match\_target}: the quantity a reproduction run must recover
to count as a match.}
\label{tab:schema-matchtarget}
\begin{center}
\small
\begin{tabular}{@{}ll@{}}
\toprule
Field & Description \\
\midrule
\texttt{config} & Experimental configuration that produces the target value. \\
\texttt{metric} & The reported metric name. \\
\texttt{value} & The target value to match. \\
\texttt{scope} & The model, dataset, or setting the value applies to. \\
\texttt{match\_bar\_kind} & Match criterion: \texttt{point\_estimate}, \texttt{direction}, \texttt{threshold}, \texttt{magnitude}, or \texttt{range}. \\
\bottomrule
\end{tabular}
\end{center}
\end{table}

\begin{table}[t]
\caption{\texttt{h100\_estimate}: the compute estimate that placed each paper
in its band.}
\label{tab:schema-h100}
\begin{center}
\footnotesize
\begin{tabular}{@{}llp{7.6cm}@{}}
\toprule
Field & Type & Description \\
\midrule
\texttt{hours} & float & Estimated H100-hours for the match target. \\
\texttt{basis\_kind} & string & \texttt{derived\_from\_config}, \texttt{paper\_reported}, \texttt{comparable\_experiment}, or \texttt{compute\_unspecified}. \\
\texttt{gpu\_count} & int or null & GPU count the paper reports. \\
\texttt{gpu\_type} & string & GPU type the paper reports. \\
\texttt{wallclock\_hours} & float or null & Reported wall-clock hours, before H100 normalization. \\
\texttt{h100\_equivalent\_multiplier} & float or null & Factor converting the reported GPU to H100-equivalent hours. \\
\texttt{basis} & string & Free-text justification for the estimate. \\
\bottomrule
\end{tabular}
\end{center}
\end{table}
